%% ma: L12-B12-0023
%% lop: 12
%% chuong: 4
%% bai: 12
%% dang: 12.12.15
%% loai: TN4
%% muc: TH
%% dap_an: "C"
%% nguon: "(NBV)-ĐỀ SỐ 19-MỖI NGÀY 1 ĐỀ THI 2025.pdf (D19, MỖI NGÀY 1 ĐỀ - Đề số 19), Phần I Câu 3"
%% kien_thuc: "Bài 11 (nguyên hàm), Bài 12 (tích phân, công thức Newton-Leibniz)"
%% trang_thai: cho-duyet
%% ly_do: "Dạng mới đề xuất 12.12.15 (hai hàm có cùng đạo hàm: nhận biết khẳng định đúng về nguyên hàm và tích phân), chờ giáo viên duyệt."
%% ghi_chu: "Đáp án gốc C đúng."
\begin{ex}
Cho các hàm số $F(x)$, $G(x)$ có đạo hàm trên $\mathbb R$. Biết rằng $F'(x)=G'(x)$, $\forall x\in\mathbb R$. Khẳng định nào sau đây là đúng?
\choice{$F(x)=G(x)$, $\forall x\in\mathbb R$}{$\displaystyle\int F(x)\,\mathrm dx=\int G(x)\,\mathrm dx+C$}{\True $\displaystyle\int_a^bF'(x)\,\mathrm dx=\int_a^bG'(x)\,\mathrm dx$}{$\displaystyle\int_a^bF(x)\,\mathrm dx=\int_a^bG(x)\,\mathrm dx$}
\loigiai{Từ $F'(x)=G'(x)$ ta suy ra $F(x)=G(x)+C_0$ ($C_0$ là hằng số), nên A, B, D nói chung không đúng (ví dụ $F(x)=x+1$, $G(x)=x$). Hai hàm $F'$ và $G'$ trùng nhau nên $\displaystyle\int_a^bF'(x)\,\mathrm dx=\int_a^bG'(x)\,\mathrm dx$. Chọn C.}
\end{ex}
